%# A dummy statement whose only job is to exercise the inline var hints.
%# Every line below is annotated with what the editor should (or should not)
%# draw beside it. The file still builds: nothing here is deliberately broken.

%- block legend
Given two integers $A$ and $B$, determine the value of $A + B$.
%- endblock

%- block input
%# BADGED. The shorthand and the long form are the same var, so both get a hint.
The input is a single line with two integers $A$ and $B$, where
$\VAR{A.min} \le A \le \VAR{A.max}$ and
$\VAR{vars.B.min} \le B \le \VAR{vars.B.max}$.

%# BADGED through the filter. rbx renders the whole expression, so the hint
%# shows what the statement will typeset, not the raw digits.
%#
%# `A.max` is a poor demonstration: 2147483647 leaves a remainder, so `sci`
%# declines to abbreviate it and hint and value agree. That is the filter
%# working, not failing.
Neither value exceeds $\VAR{A.max | sci}$.

%# These are the ones to look at. Hints: 10^5, 2x10^5, 10^18 + 7 -- rendered
%# with superscript digits, because an inlay hint cannot typeset maths.
For scale, $\VAR{DEMO.n | sci}$, $\VAR{DEMO.big | sci}$ and
$\VAR{DEMO.mod | rsci}$.

%# BADGED, unfiltered, and exact: 1000000000000000007 survives the trip whole.
%# A JSON number would have arrived as ...000, a wrong badge shown with
%# confidence, which is why rbx hands the extension display strings.
The modulus is $\VAR{DEMO.mod}$.
%- endblock

%- block output
Print one integer: the sum of $A$ and $B$.
%- endblock

%- block notes
%# NOT BADGED. `g` is the loop variable, so this line renders a different
%# value on every pass -- one hint would have to lie or name a group.
%- for g in problem.groups
Group \VAR{g.name}: up to $\VAR{g.vars.A.max}$ (shorthand: $\VAR{g.A.max}$).
%- endfor

%# NOT BADGED either: addressing a group by name is still group scope.
The samples group stops at $\VAR{problem.groups.samples.vars.A.max}$.

%# NOT BADGED: a commented-out reference is not a reference. \VAR{A.max}

%# Things worth trying by hand, none of which are in the file because each
%# would break the build:
%#   - misspell one, e.g. \VAR{A.mx}  -> the hint disappears, which is how a
%#     typo shows up: it is the only reference on the line without one.
%#   - write an expression, e.g. \VAR{A.max + 1} -> no hint, nothing is guessed.
%#   - toggle `rbx.statementVarHints`, or VS Code's own inlay-hints setting.
%- endblock
